\chapter*{Abstract}

Public Eurostat unmet-care indicators are valuable for monitoring health-access barriers, but aggregate country-year panels do not define one stable analytical population. Conclusions can change when analysts change the denominator, country universe, population weighting rule, covariate-completeness requirement, or missing-data assumption. This thesis studies that problem as target-population drift.

The thesis formalizes a diagnostic drift-stability framework. Each estimand is represented by its outcome, denominator, country universe, row set, weighting rule, missing-data assumption, and model. The framework defines a Target-Population Drift Index and a Conclusion Stability Index, but it is not a new estimator. The empirical design begins from the core finding that the primary unmet medical care indicator is observed for 608 country-years, while the baseline complete-case fixed-effects model uses 282 rows. The analysis extends this logic to a multi-outcome Eurostat benchmark covering medical and dental unmet-care indicators, population- and need-denominator measures, and feasible barrier-specific definitions.

The thesis is organized around five questions: whether monitoring conclusions depend on indicator definition; how much country, year, variable, and group coverage changes the analytical target population; whether social-vulnerability coefficients remain stable across complete-case, population-weighted, balanced-outcome, IPW, MAR-imputed, and MNAR-shifted estimands; how semi-synthetic Eurostat-style missingness affects recovery of a reference estimand; and whether a reproducible reporting protocol can classify findings as stable, target-dependent, denominator-dependent, missingness-dependent, imputation-model-dependent, or non-identifiable.

The contribution is methodological-applied. The thesis provides a public-data estimand registry, a multi-outcome monitoring benchmark, drift and stability metrics, a coefficient-stability matrix, a semi-synthetic missingness simulation, and a reproducibility package. The central object is not one poverty/social-exclusion coefficient, but the interpretation change created when public aggregate monitoring data are converted into selected analytical panels. The thesis remains aggregate, associational, and non-causal.
